\chapter{Real patterns}

\textit{Amelie showed Basil Clara.} English packs much more meaning into this terse sentence than is conveyed by the words alone. Sure, \textit{show} tells us that there is a showing situation, and \textit{Amelie} says that somebody named Amelie was involved, but we also understand that Amelie is the agent, the one doing the showing, that Basil is the recipient, and that Clara is the patient, the one being shown.

As English speakers, it's difficult for us to imagine that these could have been otherwise, but indeed they could. The same sentence in Azerbaijani, for example, would be \textit{Amelie Basilə Claranı göstərdi} (where \textit{göstərdi} means `showed'). The Cariban language Hixkaryana, spoken by roughly 500 people on the Nhamundá River in Brazil, would start the sentence with \textit{Clara} (as the patient), followed by the verb `showed', and only then would \textit{Amelie} appear (as the agent), with \textit{Basil} (as the recipient) in the final position.

But role shuffling is not the only possible language trick. A role could be assigned without anyone being named. In the imperative \textit{Show Clara to Basil}, the agent role is now yours: you are to do the showing. If you could have a role in the imperative without even being mentioned, why not in \textit{Amelie showed Basil Clara} too?

English grammar is why not. It removes all role-related ambiguity, distilling the narrative of the showing situation, with all its roles and relationships, into a pattern: \textsf{Agent + tensed verb + Recipient + Patient} plus a mere quartet of words: \textit{Amelie}, \textit{Clara}, \textit{Basil}, and \textit{showed}, plus .

It is the whole pattern that matters; the ordinal positions alone will not do. For instance, \textit{Basil} being second doesn't automatically make him the recipient. In \textit{Amelie showed Basil to Clara}, Basil retains his position but swaps roles with Clara. Even the first position, held by \textit{Amelie}, isn't a reliable indicator of agenthood. In \textit{Amelie was shown Basil}, Amelie becomes the recipient, and the agent is left unspecified.

The pattern is grammar. And this grammar has a remarkable data-compression feature: it encodes all the role-related information, so that any English speaker knows without the need for words like \textit{agent} and \textit{patient} or \textit{presenter} and \textit{audience} who is showing who to whom. Somebody who doesn't speak English, even knowing what the verb \textit{showed} means, wouldn't be able to extract the information reliably. A speaker of Azerbaijani may assume that Basil is the patient and Clara the recipient, while a Hixkaryana speaker might assume that Amelie is the patient.

In this way, the data compression feature of grammar works only for a particular social group: in this case, English speakers. And because of this, the pattern has another interesting feature: it allows us to predict how English speakers will behave. We can predict with near certainty that an attentive speaker of English will interpret the sentence with Amelie showing Clara to Basil. We can also predict with good certainty that they would express it this way.\footnote{Based on a quick and dirty search of COCA, I'd say there's about a one in five chance that they would instead say \textit{Amelie showed Clara to Basil.}}

\section{Spraying and loading}

There are more subtle patterns that linguists have noticed and puzzled over greatly but that rarely come up in popular discussions of grammar, perhaps because there's no obvious disagreement about how they are used.

There's a set of verbs, for instance, called the \textit{spray}/\textit{load} verbs. Like the verb \textit{show}, which can have A showed B C or A showed C to B, \textit{spray}/\textit{load} verbs also have two patterns that mean more or less the same thing: \textit{She sprayed the area with a fine mist} or \textit{She sprayed a fine mist over}/\textit{throughout}/\textit{into the area}. Similarly, we have \textit{They loaded the truck with hay} in alternation with \textit{They loaded hay onto the truck}. Other verbs of this type are \textit{dabbed}, \textit{heaped}, \textit{packed}, \textit{smeared}, and \textit{spread}.

Strangely enough, though, when we look at other semantically similar verbs -- those that describe distributions or accumulations, such as \textit{drench}, \textit{soak}, or \textit{flood} -- we find that they do not offer the same kind of alternation. For example, while we can say \textit{She soaked the area with water}, saying *\textit{She soaked water on the area} would be considered ungrammatical or awkward by most English speakers.

Linguists haven't been able to figure out why there are these two groups of verbs, let alone why they exhibit (or don't) such specific and consistent patterns of alternation. Despite extensive research into their syntactic and semantic properties, a comprehensive theory that accounts for the unique behavior of spray/load verbs as opposed to other semantically similar verbs remains elusive. The challenge lies not just in describing the observed patterns, but in explaining why these particular verbs group together in the first place, and why other seemingly similar verbs do not. Various hypotheses, ranging from historical linguistic developments to cognitive processing constraints, have been proposed, but none have achieved universal acceptance.

And yet, we can reliably predict that English speakers will not produce sentences like *\textit{He shoveled the garden with earth} while finding those like \textit{She smeared the bread with jam} perfectly fine.

\section{Raw arbitrary conventionality}

\citet[231]{Hurford2012a} challenges the idea that grammaticality is essentially semantic with the examples in (\ref{ex:likely-v-probable}).

\ea \label{ex:likely-v-probable}
    \ea[]{\textit{It is likely}/\textit{probable that John will leave.}}
    \ex[]{\textit{John is likely}/*\textit{probable to leave.}}
    \z
\z

The first time I read this, a few years ago, it struck me as a strong challenge. My judgments match Hurford's. But recently, it occurred to me to check, and that's how I learned about \textit{he's probable to play}.

\bigskip

The story starts with a bribery and gambling scandal in the 1946 National Football League championship game between the New York Giants and the Chicago Bears. After suspending the two Giants players indefinitely, commissioner Bert Bell realized that trust was ebbing and that something had to be done to protect the NFL from the negative influence of gambling.

There was no practical way to do away with gambling on NFL games, but one pernicious source of speculation revolved around inside information about player injuries. Bell realized that he could largely neutralize this aspect by mandating that all teams submit weekly injury reports.

In 1947, Bell introduced a requirement that the status of each injured player was to be listed as ``probable'', ``questionable'', or ``doubtful'' for the upcoming game. This information would be made available to the public, ensuring that everyone had access to the same information about player injuries.

Other major sports leagues followed suit, although the exact year of implementation varies. For example, the National Basketball Association has required teams to submit injury reports since the 1980s.

Perhaps coincidentally, I find no examples of \textit{probable to play} before the 1980s. But on Nov. 18, 1986, \textit{The Los Angeles Times} published a column by Mark Heisler that included the following example \textit{This on a day that the Raiders had him listed as \uline{probable to play} against the Browns.} Since then, in the sports media, the construction \textit{probable to play}/\textit{start}/\textit{return}/etc. has gained currency. It's still not common, mostly limited to American sports media, but it can be found from \textit{The New York Post} to the \textit{LA Times}, \textit{CBS Sports} to \textit{ESPN}, and \textit{The Toronto Sun} to the \textit{Houston Chronicle}.

In this way, \textit{he's probable to play} has come to mean something: that though an injury has been sustained, it is not currently severe enough to keep the player from the game. Consider that, although almost every starting professional athlete is \textit{likely} to play in any given game, only those who had been injured could be described as \textit{probable to play}. In major-league baseball, pitchers typically have a number of non-playing games between each of their pitching appearances, but when it's a given pitcher's turn to play, he may be \textit{likely} to play but he's probably not \textit{probable}.

Interestingly, in the National Hockey League at least, it has become common to refer to injuries simply by naming the body part: \textit{Matthews it out with an elbow}, for instance, doesn't mean that Matthews is having a night on the town with an arm joint but rather that he cannot play because he may have a cracked radial head.\footnote{I thank Michael Dipetta for this detail.}

\bigskip

As Hurford observes, one reason that \textit{likely} and \textit{probable} appear to have different structural possibilities may be that \textit{likely} comes from Old Norse, \textit{probable} from Old French. The \textit{to} of the \textit{to}-infinitival construction grammaticalized from the preposition \textit{to}, and this process was already well on its way in Old Norse. At the same time, Old French had similar expressions with \textit{à} `to' but it was just one of a number of true prepositions doing this kind of work.

While it's not clear, \textit{likely} may have already licensed the equivalent of a \textit{to}-infinitive in Old Norse, while \textit{probable} probably didn't in Old French. Thereafter, each word simply continued on its own linguistic rails until \textit{probable} picked up this new sports-based meaning.

What this illustrates is that there's nothing inherent in the structure or meaning of the word \textit{probable} that stops it from having a \textit{to}-infinitival complement in the way that adjectives like \textit{happy}, \textit{ready}, and \textit{likely} do. There's just a historical path that hadn't led that way. The injured-players list, though, provided just the right motivation because \textit{probable} is the terminology that was adopted to signify a particular status. Most of the time sports announcers leave it at that: \textit{the player is probable}. But if there's some lack of clarity about what's being discussed, it can be cleared up by appending \textit{to play}. And we have a new construction.

This shows that the forces that cause grammaticalization are not the same as those that give us the feeling of grammaticality. 

Hurford notes that a word's ``grammatical distribution does not follow completely from their meanings'' \citet[302]{Hurford2012a}. It seems to me that, in doing so, he may be confounding grammaticalization with grammaticality, what could become grammatical with what is.

He cites the example of \textit{the fifth of November}/\textit{each month}/\textit{the month in which your birthday falls} and contrasts it with \textit{November the fifth}, which works, and *\textit{each month the fifth} and *\textit{the month in which your birthday falls the fifth}, which don't.

In the Early English Books Online corpus, \uline{\textit{November the} [ordinal number]} occurs not once from the 1470s to the 1620s. Then, in the following decade, there are three examples. In the 1640s, there are 11, then 32 in the 1650s, and 67 in the 1660s. \textit{January} follows roughly the same pattern. \textit{February} has an interesting case in 1610, when we find \textit{hee died with his brother the same moneth of \uline{february the twentith day}}, which seems to be an intermediate form, and \textit{March} has a similar example in 1592 \textit{dated from Groyne in march the thirtith day}.\footnote{There appear to be scattered instances from the 1550s, but these are spurious hits, such as \textit{february the seconde yeare of her highnes most noble raygne}.}

Perhaps the earliest example of the [month] \textit{the} [date] format is from 1616.  

\begin{quote}
    \textit{therefore the best ground to sow hempe on, is the richest of all mixt earths, whose mould is driest, loosest, blackest, and quickliest ripe, with little earing, as namely with two ardors at the most, which vvould be in \uline{october the first}, and the last in march, which is the best and most conuenient time for sowing: }
\end{quote}

Notice the use of \textit{in} with \textit{october the first}, putting us in the month of October, where today we would likely say \textit{on October the first}, putting us on a date.

At any rate, there appear to be no genuine examples before the seventeenth century, and the format seems to be growing through the middle of the century, and establish before the following century rolls around.

Interestingly, Chaucer had \textit{the eightetethe} `18th' \textit{day of April} near the end of the fourteenth century, but \textit{the} [date] \textit{of} [month] format only really becomes established in the sixteenth century. Prior to that, it was also common to refer to particular landmark days such as feast days or the ides `middle' of the month and then to calculate forward or back. Date formats seemed to be in flux.

Hurford claims that the \textit{November the fifth} construction imposes a semantically unmotivated restriction to just the proper names for months. It's unlikely that a new semantic distinction was being made with the \textit{November the fifth} format. It was more likely some kind of social or register distinction. The rise of bureaucracy and record-keeping could have contributed. Written documents benefited from brevity and standardization. The format could have initially been motivated by these bureaucratic reasons, subsequently becoming seen as more formal or ``educated'', particularly compared to the less standardized expressions of date, such as \textit{Five days before the Kalends} `the first' \textit{of May}.

The counterexamples *\textit{each month the fifth} and *\textit{the month in which your birthday falls the fifth}, are notably the kind of thing which does not seem to suit bureaucratic standardization.

Part of Hurford's concern is that, for computational linguists, a large number of such rules that apply to only small groups of words such as the twelve month names is computationally inefficient. But the recent success with having Large Language Models, such as Chat GPT, at being grammatical, not only at the general grammatical facts, but also at particular micro-constructions such as \textit{November the fifth}, suggests that perhaps we simply need to bite that computational bullet.

\bigskip

Chris Anderson's article titled ``The long tail'' appeared in \textit{Wired} magazine in October 2004. Anderson developed the concept into a book, published in 2006. The term \textsc{long tail} refers to the strategy of businesses that sell a large number of unique items in relatively small quantities, in contrast to selling a small number of items in large quantities. This concept is particularly relevant in the context of digital goods and online markets, where the cost of inventory storage and distribution is significantly lower than in traditional retail settings.

The popularity of the "long tail" concept has led to its application in various fields, including statistics, business models, and information theory, reflecting its broad impact on understanding and leveraging the distribution of products, services, and information in the digital age. Hurford's concerns about the micro-grammar of words like \textit{November} are similar to the problems that merchants faced before the internet. Computational linguistics didn't have the massive corpora or the sheer amount of compute needed to learn the micro-grammars, but 



There seems to be an ability in these models to do few-shot learning, as long as some element of the situation is novel \citep{Liu2019}.

% TODO: write a functionalist response to Hurford's challenge in Brett's voice — earlier draft here was LLM-generated and has been removed.

% consider a simulation like that in the code.py run in https://trinket.io/embed/python3/a5bd54189b
% this simulation doesn't work well

